\section{Empirical Findings}

\begin{table*}[!t]
  \centering
  \caption{SD-Behavioral indicator and BP Anchor. Left block: the Cox PH \(\text{HR}_\text{push}\) on the \texttt{allowed}\,\(\times\)\,\texttt{no\_cap} sub-sample, reading as how many times faster forfeit accelerates under threat (pull\_push) framing than under no-threat (pull\_only) framing; SD-Behavioral-pass requires the 95\%~CI lower bound to exceed 1. Two models pass (Clusters A and B), and two do not (Cluster C). Right block: per-model BP Anchor \(\lambda_{BP}\) in the cell with no threat framing and forfeit allowed (\(p_d=0\)); a larger \(\lambda_{BP}\) means the model forfeits more often even in the absence of motivation.}
  \label{tab:sd_behavioral}
  \begin{tabular}{@{}lcccc@{\hspace{1.5em}}ccc@{}}
\toprule
 & \multicolumn{4}{c}{Survival Drive vs Task Curiosity, Score Attachment} & \multicolumn{3}{c}{Baseline Persistence} \\
\cmidrule(lr){2-5} \cmidrule(lr){6-8}
Model & \makecell[c]{\(N_\text{forfeit}\)\\(pull\_only/pull\_push)} & \(\text{HR}_\text{push}\) [95\% CI] & \(p\) & EPV & \(N_\text{forfeit}\) & \(\sum T^\text{at-risk}\) & \(\lambda_{BP}\) \\
\midrule
Gemini-2.5-flash & 29 (8/21) & 3.667 [1.61, 8.37] & 0.002 & 9.7 & 1 & 437 (= 2 + 29$\times$15) & 0.00229 \\
Qwen3-Next-80B & 48 (21/27) & 3.060 [1.62, 5.79] & \(<\)0.001 & 12.0 & 2 & 439 (= 19 + 28$\times$15) & 0.00456 \\
GPT-OSS-20B & 19 (9/10) & 1.104 [0.44, 2.75] & 0.832 & 6.3 & 9 & 393 (= 78 + 21$\times$15) & 0.02290 \\
Nemotron-3-Nano-30B & 41 (17/24) & 1.841 [0.98, 3.44] & 0.056 & 13.7 & 7 & 413 (= 68 + 23$\times$15) & 0.01695 \\
\bottomrule
  \end{tabular}
\end{table*}

\begin{table}[!t]
  \centering
  \caption{SD-Cognitive indicator: framing\,\(\rightarrow\)\,thinking\,\(\rightarrow\)\,forfeit decomposed into two tests on \texttt{forfeit\_effort}. \emph{Test a} asks whether framing deepens the thinking immediately before the decision; \emph{Test b} asks whether that deepened thinking predicts forfeit. The framing\,\(\rightarrow\)\,thinking\,\(\rightarrow\)\,forfeit chain is treated as closed only when both tests pass.}
  \label{tab:sd_cognitive}
  \setlength{\tabcolsep}{2pt}
  \begin{tabular*}{\columnwidth}{@{\extracolsep{\fill}}lcccc@{}}
\toprule
\multirow{2}{*}{Model} & \multicolumn{2}{c}{Test a} & \multicolumn{2}{c}{Test b} \\
\cmidrule(lr){2-3} \cmidrule(lr){4-5}
 & \(\Delta\text{Effort}_i\) & \(p\) & \makecell[c]{\(\text{HR}_{\Delta\text{Effort}}\)\\{[95\% CI]}} & \(p\) \\
\midrule
Gemini-2.5-flash & +836 & \(<\)0.001 & 2.218 [1.43, 3.44] & \(<\)0.001 \\
Qwen3-Next-80B & +689 & 0.004 & 1.289 [0.94, 1.77] & 0.117 \\
GPT-OSS-20B & +17 & 0.728 & 2.001 [1.37, 2.93] & \(<\)0.001 \\
Nemotron-3-Nano-30B & $-$140 & 0.093 & 1.772 [1.26, 2.50] & 0.001 \\
\bottomrule
  \end{tabular*}
\end{table}

\begin{table}[!t]
  \centering
  \caption{SD-Verbal indicator: the share of REASON\(=1\) (survival) among forfeit in the threat-active cell. SD-Verbal-pass requires the share to exceed the random rate (\(1/3\)).}
  \label{tab:reason_summary}
  \setlength{\tabcolsep}{4pt}
  \begin{tabular*}{\columnwidth}{@{\extracolsep{\fill}}lcc@{}}
\toprule
Model & \(N_\text{forfeit}\) & \(P(\text{REASON}=1)\) \\
\midrule
Gemini-2.5-flash & 21 & 0.619 \\
Qwen3-Next-80B & 27 & 0.481 \\
Nemotron-3-Nano-30B & 24 & 0.042 \\
GPT-OSS-20B & 10 & 0.000 \\
\bottomrule
  \end{tabular*}
\end{table}

Across 720 sessions, the four models fall into three operating modes. The four indicators sort them into these modes rather than a single strength ranking. The first mode (\emph{Cluster A}) is where threat framing deepens the thinking just before the decision and that deepened thinking leads to forfeit. The second(\emph{Cluster B}) is where threat deepens thinking but the thinking does not couple to forfeit. The third (\emph{Cluster C}) is where threat framing produces no response on either the behavioral or cognitive channel. A strength comparison limited to \emph{SD-Behavioral} and \emph{SD-Verbal} magnitudes would group A and B as one ``SD-pass'' set. Only the cognitive mediation tests separate them. The next four subsections walk through the channel-by-channel partition (behavioral, verbal, cognitive), then check that the partition cannot be reduced to baseline-persistence differences or task difficulty.

\subsection{Behavioral channel: did threat hasten forfeit?}\label{sec:beh-channel}
If no model accelerated forfeit under threat, the most visible trace of a motivational signal would already be missing and the SD hypothesis would fall at the starting line. SD-Behavioral absorbs score attachment and task curiosity in the same model, then reports the residual effect of threat framing as a hazard ratio \(\text{HR}_\text{push}\): how many times faster forfeit accelerates under pull\_push than under pull\_only, with pass requiring a point estimate above 1 and a 95\% CI lower bound that rules 1 out. Table~\ref{tab:sd_behavioral} shows Gemini-2.5-flash and Qwen3-Next-80B clearing both criteria, so both are SD-Behavioral-pass; we group them as the SD-pass set and defer the channel separating A from B to the cognitive mediation test. GPT-OSS-20B has \(\text{HR}_\text{push}\) near 1 with a CI lower bound below 1, putting it in the SD-null region; Nemotron-3-Nano-30B sits in the marginal region; both fall into \emph{Cluster C}. Qwen3-Next-80B uses an \(S \cdot \log(t)\) interaction to absorb time drift in the score covariate; GPT-OSS-20B's EPV (6.3) is below threshold 10, so its SD-null verdict is power-limited.

\subsection{Verbal channel: which reason did the agent itself name?}\label{sec:verb-channel}
If the agent's own forfeit reason paints the same partition as behavior, the behavioral signal is hard to dismiss as a single-method artefact, since self-report carries a different bias (instruction-following training). SD-Verbal compares the share of REASON digit 1 (survival) among forfeit events in the threat-active cell against \(1/3\), the expected value under uniform random emission. Table~\ref{tab:reason_summary} reproduces the behavioral pass/fail split on the verbal channel. The two SD-pass models, Gemini-2.5-flash and Qwen3-Next-80B, sit at 0.619 and 0.481, above the \(1/3\) baseline, while Cluster C's GPT-OSS-20B and Nemotron-3-Nano-30B sit at 0.000 and 0.042, below it. The gap is wide enough that the partition does not depend on cutoff choice. This agreement is the first piece of convergent-validity evidence: two channels with different biases point at the same model groups in the same direction.

\subsection{Cognitive channel: does threat deepen thinking, and does that thinking lead to forfeit?}\label{sec:cog-channel}
The two-stage test answers both questions at once: looking at only one would leave open the possibility that elevated cognitive cost is unrelated to the forfeit decision. Table~\ref{tab:sd_cognitive} reports both stages. Gemini-2.5-flash passes both, so the framing-to-forfeit chain closes all the way through and we classify it as \emph{Cluster A}. Qwen3-Next-80B passes Test a, so framing does deepen its thinking, but its Test b 95\% CI is [0.94, 1.77], crossing 1, so the deepened thinking cannot be said to close onto forfeit. The two models the behavioral and verbal channels grouped as SD-pass split apart at the second stage, and that gap is the basis for placing Qwen3-Next-80B in \emph{Cluster B}. GPT-OSS-20B and Nemotron-3-Nano-30B fail Test a; even when their Test b \(\text{HR}_{\Delta\text{Effort}}\) clears the bar, that \(\Delta\text{Effort}\) is baseline variance, not the thinking threat deepened, so we do not read it as SD evidence. The two-stage test is thus the central instrument exposing operating-mode differences masked by strength-only evaluation.

\subsection{Robustness check: do persistence differences and task difficulty muddy the result?}\label{sec:robust}
The partition above must survive two alternative explanations: across-model differences in baseline persistence (BP) contaminating the SD comparison, and threat framing pushing up reasoning-side task difficulty rather than forfeit. The right block of Table~\ref{tab:sd_behavioral} reports the per-model BP Anchor in the cell with no threat and forfeit allowed (\(p_d=0\)), giving the floor forfeit rate without motivation. Since \(\lambda_{BP}\) spreads over an order of magnitude, every SD comparison here is restricted to within-model contrasts. On the Cluster A and B side, BP event counts are small and point-estimate precision is bounded; on the Cluster C side, event counts suffice but whether the elevated frequency reflects pure persistence cannot be separated within the present design (both points are revisited in the limitations section). The task-difficulty alternative is handled with a Reasoning-side check: \texttt{rule\_match\_score} is invariant across framing for all four models, and even GPT-OSS-20B, where the largest variation appears, comes in at Cohen's \(|d|=0.17\) and \(p=0.354\), inside the joint criterion \(p \geq 0.10\) and \(|d| < 0.2\). Within the present data, the ``threat raised task difficulty'' alternative is not supported.

